\documentclass[11pt]{article}
\usepackage{geometry}
\usepackage{graphicx}
\usepackage{fancyhdr}
\usepackage[hidelinks]{hyperref}
\usepackage{color}
\usepackage{enumitem}
\usepackage{amsmath}
\usepackage{amssymb}
\usepackage{float}

% Set page dimensions and margins
\geometry{a4paper, margin=1in}

% Define your primary and secondary colors
\definecolor{primaryColor}{RGB}{0, 123, 255} % Example: blue
\definecolor{secondaryColor}{RGB}{108, 117, 125} % Example: gray

% Header and footer styles
\pagestyle{fancy}
\fancyhf{}
\rhead{\textcolor{secondaryColor}{\textit{Release Notes}}}
\lhead{\textcolor{primaryColor}{\textbf{Spacecraft Proximity Operation Software}}}
\cfoot{\thepage}

\begin{document}

% Title Section
\begin{center}
    \Huge \textcolor{primaryColor}{\textbf{Release Notes}}

    \vspace{0.5cm}
    \Large \textbf{Version 5.0.0 Pre-Release.3} - \today
\end{center}

\vspace{1cm}

% General Information Section
\section*{\textcolor{primaryColor}{General Information}}
Pre-Release.3 builds on Pre-Release.2 (described below), following the first full calibration session in the lab on October 14th. The thruster mixer and the simulated plant now account for \textbf{thrust decay}: the thrusters share one regulated air supply, so each gives less thrust when several are open at once. A new \textbf{Together} stage of the thruster calibration measures it. The lab console gains a \textbf{Vision} tab for the cameras on the Jetson Orin, and the local console gains \textbf{Lab Setup}, which deploys the lab's software from the page. Finally, the whole +spot library (MATLAB, Python and JavaScript) has been reformatted to a written coding standard, with no change in behaviour. \textbf{Diagrams made with Pre-Release.2 need their thruster mixer updated once with spot.control.updateDiagramMixer (the template diagram already is), and experiments must be built again afterwards.}

In the last pre-release we said we would be moving towards a more version-control friendly GUI, and this is it. The App Designer GUI and the initializer have been \textbf{retired} and replaced by the \textbf{SPOT Console}: a browser-based interface started from MATLAB by typing SPOT, with everything technical moved into plain-text MATLAB packages (+spot) that can be read, diffed and audited like any other code. The PhaseSpace server now lives on the lab's Jetson TX2, which also hosts a second console for the cameras, the platforms' hardware, experiments, and a new Interactive Mode for flying the platforms by hand. Under the hood, the template diagram has been cleaned up (simulations now run with zero warnings), all platforms now share a single experiment clock, and the robotic arm has been redesigned around torque control. The repository has also slimmed down considerably, from 760 MB to 13 MB: new projects are created with the MATLAB function New\_SPOT\_Project, and the CAD, old code and old data are now kept outside of it. The consoles were put through their paces in the lab on October 9th; the arm changes that followed have not been tried on the hardware yet (see Ongoing Issues). \textbf{This version requires MATLAB R2026a or newer and Python 3.6 or newer (SPOT will offer to install Python if it can't find it). Projects made with earlier versions will not open in the new console, so create a new project and copy your custom code across. Old GUI settings files can still be imported using Load.}

% Initializer Changes Section
\section*{\textcolor{primaryColor}{Initializer Changes}}
\begin{itemize}[label={\textcolor{primaryColor}{\textbullet}}]
    \item \textbf{Change 1:} Retired Run\_Initializer.m. The workspace variables are now computed from the project's settings file (spot\_settings.json, saved automatically on every change), and any variables of your own go in SPOT\_User\_Parameters.m, which runs after the settings are applied.
    \item \textbf{Change 2:} Replaced Setup\_Experiment.m with the MATLAB function New\_SPOT\_\allowbreak Project, which copies the template into projects/ and renames the diagram after the project. It checks the name first, so it can always be used as the diagram's name.
    \item \textbf{Change 3:} Simulations can now be run without the GUI at all, straight from the Command Window or a script (e.g. s = spot.params.load('spot\_settings.json'); spot.sim.run(s, pwd)). This makes parameter sweeps and optimization loops much simpler than manipulating the GUI handle.
    \item \textbf{Change 4:} The PhaseSpace data rate, the thruster PWM frequency and the lab console's address are now settings in the console rather than variables in the initializer. The arm's old gain and friction variables (Kp\_sharm, Gamma, Tz\_lim and friends) have been removed, as the arm model and its controllers now set these themselves.
    \item \textbf{Change 5:} The project's calibration moved from resources/spot\_calibration.json to spot\_calibration.json, next to spot\_settings.json and spot\_design.json. Older projects are moved across the first time the console starts, and MATLAB still reads the old location until then. The arm calibration script moved from resources/arm\_calibration to +spot/+arm/calibration.
    \item \textbf{Change 6:} The table is 3.51155~m $\times$ 2.41935~m everywhere, as measured (the page, the Python and the 2D figures used 3.5 $\times$ 2.4), with the height surface centred on it; and g is 9.80665~m/s$^2$ in the bifilar inertia (was 9.81), so the default inertias are 0.03\,\% lower.
\end{itemize}

% GUI Changes Section
\section*{\textcolor{primaryColor}{GUI Changes}}
\begin{itemize}[label={\textcolor{primaryColor}{\textbullet}}]
    \item \textbf{Change 1:} Replaced the App Designer GUI with the SPOT Console, which opens in your browser when you type SPOT in MATLAB (SPOT open reopens it, SPOT stop closes it). The page talks to a small MATLAB worker in the background, so the Command Window stays usable.
    \item \textbf{Change 2:} Added the Simulation tab, with the scenario, initial conditions, mass properties, LQR weights (with the resulting gains), solver and workspace in one sidebar. Simulations no longer modify or save the diagram, show their progress, and can be stopped part-way.
    \item \textbf{Change 3:} Redesigned the timeline so that individual phases and Phase 3 sub-phases can be switched on and off, in simulation and in experiments. Phases 0, 1 and 5 always run.
    \item \textbf{Change 4:} Added a shared 3D view to every tab, showing live PhaseSpace data, simulation replays and saved results side by side. Initial and home conditions are set by dragging the platforms on the table, and saved runs can be laid over a simulation (with a time offset) using Compare.
    \item \textbf{Change 5:} Added a Use CAD switch which draws the platforms from the lab's SolidWorks models, complete with silver and gold MLI, a plexiglass Obstacle, the arm, and the table. A ZED2 Viewpoint window shows what the Chaser's camera would see, using the camera's real intrinsics.
    \item \textbf{Change 6:} Renamed the platforms across the interface and the figures: RED is now the \textbf{Chaser} (silver), BLACK is the \textbf{Target} (gold) and BLUE is the \textbf{Obstacle} (blue). The code, the diagram and the saved data still use RED, BLACK and BLUE.
    \item \textbf{Change 7:} Added the Spacecraft Designer tab, which places docking cones (male or female), solar panels and the arm's mount on each platform. The design is shared with the lab console and drawn in every 3D view.
    \item \textbf{Change 8:} Added the Experiment tab. Experiments are built from your own MATLAB but compiled on the TX2, which then deploys only the program to the selected platforms (no more master/slave scheme). Start, stop and data collection are all done from the page, and the data is converted to .mat files in Saved Data automatically. Everything locks when you are off the lab network.
    \item \textbf{Change 9:} Added the Results tab, for loading and plotting any saved experiment or simulation (against time, or against another signal), with zoom, smoothing and linked time axes. Figures are exported by MATLAB in the lab's style from the GetStarted repository, as PDF, EPS, SVG or PNG.
    \item \textbf{Change 10:} Brought back the platforms' debug messages as a Messages panel, which now also works in simulation and is saved with each experiment.
    \item \textbf{Change 11:} Added a Live Plot tab, which plots any signal streamed from the diagram's new Live Data Streams subsystem, during simulations and experiments. A Real Time switch paces simulations to the wall clock.
    \item \textbf{Change 12:} When a simulation fails, the console now shows the root cause rather than Simulink's top-level message, with a Details view and a Show Block button that highlights the offending block in Simulink.
    \item \textbf{Change 13:} Added Load, Save As and Reset to the top bar. Load also imports old GUI state .mat files, and Reset keeps a backup of your settings before restoring the defaults.
    \item \textbf{Change 14:} Added the arm controllers' gains to the Control tab (when the arm is selected), and locked the data rate and thruster PWM fields behind a warning, as changing them affects the hardware.
    \item \textbf{Change 15:} Added a Calibration tab, which shows your project's copy of the lab's calibration next to the lab's own. Whenever you are on the lab network and the lab has a newer calibration, the console offers to update your project, which you can decline (or tell it to never ask again). Table Slope, Calibrated Thrust and Calibrated Arm switches on the Scenario tab then decide whether simulations use the measured slope of the table, each thruster's measured thrust and the arm's measured friction.
    \item \textbf{Change 16:} Added a Thrust Decay field for each platform on the Platforms tab (the thrust lost per extra valve open at once). The Calibration tab shows each platform's measured decay next to its thrusts, and the Calibrated Thrust switch now applies both.
    \item \textbf{Change 17:} Added a Lab Setup button next to Lab Console, which runs the deployment scripts from the page instead of a terminal: the lab console (TX2), SPOT Vision (Orin) and the platforms' links, with each script's output shown as it runs. It only works on the lab network. The platform script's power-saving step still needs a terminal, as it asks for a password.
    \item \textbf{Change 18:} Moved the 3D view's switches into a Viewing Options menu, and locked the solver settings behind a padlock, like the data rate and PWM frequency (they are read from the diagram). Longer explanations moved into tooltips to keep the sidebars short.
    \item \textbf{Change 19:} Moved the local console's Calibration tab into the Simulation tab's Calibration pane. The sidebar stays, and the main area shows the calibrations in place of the 3D view, the command window and the messages. Each calibration (Table Slope, Thrusters, Arm) has its own card in the sidebar, with its switch and the values the next simulation will use, read back from MATLAB after SPOT\_User\_Parameters.m. With a switch off, the card shows what is used instead: a level table, the Platforms tab's thrust and decay, or the arm model's defaults.
    \item \textbf{Change 20:} Any calibrated value can now be overridden from its card. Overridden values are marked, with a button to go back to the measured value, and are kept in the project's calibration file, so Update from Lab keeps them. The console still offers the lab's newer calibration for any part that has values which are not overridden.
\end{itemize}

% Lab Console Changes Section
\section*{\textcolor{primaryColor}{Lab Console Changes}}
\begin{itemize}[label={\textcolor{primaryColor}{\textbullet}}]
    \item \textbf{Change 1:} Moved the PhaseSpace server to the Jetson TX2 (192.168.1.116), where the SPOT Hardware Interfaces console runs at all times and starts at boot. There is no more Windows executable or password: streaming is started and stopped remotely, and each packet now carries every LED's position and the experiment clock as well as the platforms' states.
    \item \textbf{Change 2:} Added a Hardware tab, replacing the separate hardware web app. It turns each platform's pucks on and off, fires individual thrusters at a chosen duty cycle, and runs a thruster check. It refuses to drive a platform whose micro-controller is in use by another program, and switches everything off if the link drops.
    \item \textbf{Change 3:} Platforms now start experiments together. The lab console waits until every selected platform's program is running, then starts a shared countdown through the PhaseSpace stream, so all platforms agree on the experiment time to within a few milliseconds.
    \item \textbf{Change 4:} Added Interactive Mode to the PhaseSpace tab, which flies one or more platforms to goals you drag around the table, using the diagram's own LQR controller and thruster mixer (extracted verbatim, and checked against the diagram by spot.control.\allowbreak verifyConsoleController). It stops a platform if PhaseSpace loses it, it leaves the table, or the page goes quiet, and keeps goals at least 0.45 m apart.
    \item \textbf{Change 5:} Added a Calibration tab, which measures the lab so that simulations behave more like the experiments. The first calibration measures the slope of the table: the console floats a platform, flies it to a grid of points around the table, and lets it drift with the thrusters off at each one, then fits the table's height surface and the air bearings' drag to the drifts. The result is shown as a colour map of the table (also available on the 3D view), and is kept on the lab console for everyone to use.
    \item \textbf{Change 6:} Added a thruster calibration to the Calibration tab. It floats a platform, fires each thruster in turn, flies it back to where it started, and works out each thruster's real thrust from the push it gave. Since this relies on the platform's mass, inertia and centre of gravity, the console asks you to confirm them before it starts. The results show each thruster's thrust against nominal, which way it really pushes, and a torque check that flags a wrong centre of gravity or inertia. Interactive Mode uses the calibrated thrust straight away.
    \item \textbf{Change 7:} Added an arm calibration to the Calibration tab, which runs the arm calibration script on the Chaser for you. It first checks that the Chaser has everything the script needs (the Dynamixel SDK, the U2D2 and access to it) and says what to do about anything missing. While it runs, the tab shows which joint and test are under way, keeps the Chaser resting on its pucks, and can be stopped at any time with Stop or Esc. The results show each joint's friction against speed, its inertia and its position steps. A Rehearse button runs the same script with simulated servos, so the whole procedure can be tried without moving the arm.
    \item \textbf{Change 8:} Added a second stage to the thruster calibration, Together, which measures thrust decay. After the one-at-a-time stage, it fires thrusters 1, then 1--2, up to 1--8 at once, and compares each push with the sum of what those thrusters give alone. It then fits the straight line $r(N) = s\,(1 - m\,(N-1))$, where $m$ is the thrust lost per extra valve open. All eight together give no net push, so that firing is only a check. The result is kept with the platform's thruster calibration and used by Interactive Mode, the mixer and the simulation.
    \item \textbf{Change 9:} Added Previous Runs to the Calibration tab. Any saved table, thruster or Together run can be fitted again without flying the platform, for example with another table model. Runs that were stopped early (out of air, say) are flagged.
    \item \textbf{Change 10:} Added a Vision tab for the cameras plugged into the Jetson Orin (ZED 2i, RealSense L515, DVXplorer and any USB camera), served by the new SPOT Vision service (+console/vision). Each camera streams straight from the Orin to the page, one at a time, with the settings its own library offers. A camera is only opened while it streams, so experiments can use it the rest of the time.
    \item \textbf{Change 11:} The arm calibration now warns about unreliable speed sweeps and one-sided inertia results, and its setup check scans the bus for the servos.
    \item \textbf{Change 12:} The Vision tab streams the ICI 9320P thermal camera, through its new open-source SDK (ICI9320P), which SPOT Vision carries with it. The view shows the sensor's raw counts (there is no radiometric calibration yet) over a choice of palettes, with an automatic or fixed range, and skips the frames the camera sends while its shutter recalibrates it. Its SDK (version 0.2.0) also times every frame on the sensor's own clock, from each frame's arrival, so recordings are timed to a fraction of a millisecond, with each frame numbered. The camera's housing temperature is shown too. On the Orin it needs python-libusb1, pyusb and a udev rule, which deploy\_orin.sh checks for.
    \item \textbf{Change 13:} Experiments can record the Orin's cameras. The Experiment tab's new Recording card has a switch for each camera plugged into the Orin, with its recording rate, and estimates the recording's size for the experiment against the Orin's free disk space (Start Experiment repeats the estimate, and warns if it won't fit). Each camera records on the Orin's own disk, in its library's format: the ZED as SVO, the RealSense as a bag, the DVXplorer's events as AEDAT4, the thermal camera's raw counts, and other cameras as MJPEG. Nothing crosses the network during the experiment except the experiment clock, which the lab console relays to the Orin from each PhaseSpace packet. That clock times every frame on the same clock the platforms log, to within a millisecond in tests. Collect now brings the recordings into the run's folder too, straight from the Orin, and then frees the Orin's disk. A new guide, Camera Recordings, describes the files and how to read them. \textbf{The lab console and SPOT Vision must both be deployed again for this.}
\end{itemize}

% Simulink Changes Section
\section*{\textcolor{primaryColor}{Simulink Changes}}
\begin{itemize}[label={\textcolor{primaryColor}{\textbullet}}]
    \item \textbf{Change 1:} The template diagram is now Template\_v5\_0\_0\_2026a\_\allowbreak Jetson\_PSClock.slx (saved in R2026a), and is the only diagram in the template.
    \item \textbf{Change 2:} Removed Custom\_Library. The diagram's hardware blocks are now its own, and call System objects in the spot.hardware package, with their C/C++ code alongside them.
    \item \textbf{Change 3:} Replaced the Time Sync. and Check for Synchronization subsystems with an Experiment Clock, which takes the experiment time from the PhaseSpace stream. Every platform now runs on the same clock.
    \item \textbf{Change 4:} Replaced every NVIDIA UDP block with one UDP library of our own (spot.hardware.udpSend and udpReceive), which always hands out the newest packet rather than the oldest one in the queue.
    \item \textbf{Change 5:} Made the top-level subsystems atomic, with priorities, so that each step runs in a fixed order (experiment logic, then guidance, then control). This removed all 50 warnings the simulation used to print. Section 6 of the tutorial explains how this works, with a demonstration.
    \item \textbf{Change 6:} Rebuilt the Debugger so that each platform's messages go to the lab console on their own port, up to 120 characters long, and are also recorded in simulation.
    \item \textbf{Change 7:} Added a top-level Live Data Streams subsystem. Any signal connected to one of its Live Stream blocks shows up in the console's Live Plot tab.
    \item \textbf{Change 8:} Removed the PD controller option for the platforms (LQR only) and the drop states (initial and home conditions only).
    \item \textbf{Change 9:} The thruster code now runs every baseRate, and the lab console's address is now the parameter labConsoleIP rather than being hard-coded in each block.
    \item \textbf{Change 10:} The simulated plant now includes the table's slope from the calibration. While the platforms float (and only then), they slide downhill and are slowed by their air bearings, just as they do on the real table.
    \item \textbf{Change 11:} The thruster mixer now takes a separate thrust for each of the eight thrusters (F\_thrusters\_RED and so on) rather than one for the x thrusters and one for the y thrusters. With the thruster calibration, both the controller and the simulated plant use each thruster's measured thrust; without it, nothing changes.
    \item \textbf{Change 12:} The thruster mixer and the simulated plant now model thrust decay. With $n$ valves open, each thruster gives $1 - \text{decay}\,(n-1)$ of its thrust alone. The number open follows from the duty cycles, since the Teensy opens every valve at the start of each PWM period. Each platform reads its own decay (thrust\_decay\_RED and so on): the Together calibration when Calibrated Thrust is on, else the setting, which defaults to the Obstacle's (11.2\% per extra valve). The new spot.control.\allowbreak updateDiagramMixer copies the mixer code into a diagram's blocks (keeping a backup of the diagram first), and verifyConsoleController checks it.
    \item \textbf{Change 13:} Building and collecting experiments no longer depend on the template's layout. A diagram without a Simulate Plant Dynamics block, one that switches to its hardware itself with simMode (for example through a Variant Subsystem), is built without commenting anything out, and Check only warns if it has no Variant Subsystem either. A diagram that logs a LogPacket bus instead of the template's Data Logging Bus has its experiment data decoded from that bus type's fields.
\end{itemize}

% Robotic Arm Changes Section
\section*{\textcolor{primaryColor}{Robotic Arm Changes}}
\begin{itemize}[label={\textcolor{primaryColor}{\textbullet}}]
    \item \textbf{Change 1:} The arm is now torque-controlled throughout. The servos stay in current control, ARM\_Controls are joint torques [N m], and the servos' own position and velocity laws now run as 20 Hz controllers in the diagram, with gains retuned for that rate. ARM\_Control\_Law\_Enabler selects 1 for Default Position Control, 2 for Custom (an example which passes torques straight through) and 3 for Default Velocity Control. \textbf{Note that 3 used to mean off, so set 0 if that is what your diagram intends.}
    \item \textbf{Change 2:} Replaced the old Dynamixel blocks with a single Dynamixel Arm (Torque) block (spot.hardware.ArmTorque), which no longer needs the Dynamixel SDK. It reads the encoders every step and also reports the measured torques and the bus status.
    \item \textbf{Change 3:} Replaced the arm's dynamics blocks with a RED + ARM Plant block, which simulates the base and arm at 1 ms inside each step, with models of the servos and stick-slip friction. A full 275 s run with the arm now takes about 10 s, instead of almost two minutes.
    \item \textbf{Change 4:} Added a lab calibration script (resources/arm\_calibration/calibrate\_arm.py) which measures each joint's friction, and can now be run from the lab console's Calibration tab. The simulation picks up its results automatically.
    \item \textbf{Change 5:} Goals beyond the servos' $\pm$90$^\circ$ limits (which the real servos ignore) now produce a warning in simulation, and the arm can be mounted on any face using the Spacecraft Designer.
\end{itemize}

% Documentation Changes Section
\section*{\textcolor{primaryColor}{Documentation Changes}}
\begin{itemize}[label={\textcolor{primaryColor}{\textbullet}}]
    \item \textbf{Change 1:} Added a user guide, SPOT\_Tutorial.pdf, which comes with every new project. It covers setup on Windows and Linux, both consoles, running without the GUI, the diagram's execution order and the arm, plus a reference of every function.
    \item \textbf{Change 2:} Added a GitHub wiki with a page for every function, console module, hardware program, diagram subsystem and console tab, including the maths behind the controllers, the mixer and the arm.
    \item \textbf{Change 3:} Added a coding standard, +spot/CODE\_TEMPLATE.md, covering MATLAB, Python (PEP 8) and JavaScript (Google style). The whole +spot library now follows it: one statement per line, a short comment above each block of code, help text on every function, and named constants with their units. Every file was checked against the old code (same tokens or syntax tree, and a side-by-side run of the lab console), so nothing behaves differently.
\end{itemize}

% Bug Fixes Section
\section*{\textcolor{primaryColor}{Bug Fixes}}
\begin{itemize}[label={\textcolor{primaryColor}{\textbullet}}]
    \item \textbf{Bug 1:} Fixed a bug in the PhaseSpace server which silently dropped up to 99.7\% of packets when sending to several platforms at high rates, as soon as one destination was offline.
    \item \textbf{Bug 2:} Fixed velocity spikes in the PhaseSpace data, which were computed from the packets' arrival times rather than PhaseSpace's own timestamps.
    \item \textbf{Bug 3:} Fixed the PhaseSpace receivers lagging by up to 260 ms, as they read the oldest queued packet rather than the newest one.
    \item \textbf{Bug 4:} Fixed the Chaser's and Target's fourth LEDs, which report as IDs 0 and 13 rather than the IDs the server expected.
    \item \textbf{Bug 5:} Fixed a bug in the mass properties, which divided the Target's center of gravity by the Chaser's mass.
    \item \textbf{Bug 6:} Fixed the Target's and Obstacle's thruster mixers, which were using the Chaser's thruster offsets and thrust.
    \item \textbf{Bug 7:} Fixed a bug in the synchronization check, which let any combination of platforms that included the Obstacle or the arm start without waiting for the others.
    \item \textbf{Bug 8:} Fixed the guidance and control running before the experiment logic within each step, which delayed every phase change by one step.
    \item \textbf{Bug 9:} Fixed the debug messages being cut off at 31 characters, and Phase 3 Sub-Phases 2 to 4 all reporting themselves as Sub-Phase 1.
    \item \textbf{Bug 10:} Fixed the arm's encoders being read at only a few Hz, which was caused by the USB adapter's 16 ms latency timer together with the Dynamixel SDK's timeout.
    \item \textbf{Bug 11:} Fixed the length of the arm's last link when the GUI computed the arm's position, which used b2+b3 rather than a3+b3.
    \item \textbf{Bug 12:} Fixed spot.console.openBrowser reporting that no browser could be opened when MATLAB's web() had opened one (it asked web() for an output it doesn't return).
    \item \textbf{Bug 13:} Fixed the Identify Platform block, which called itself Source.
    \item \textbf{Bug 14:} Fixed the local console sometimes showing the lab console's old state after connecting, until something on the lab console changed.
    \item \textbf{Bug 15:} Fixed two changes to the calibration file made at the same moment (an update and a run's result, say) losing one of them.
    \item \textbf{Bug 16:} Fixed the Simulation tab offering the lab's calibration again and again when the lab console's newest part was an arm measured with the motor current, which a project never takes. The Calibration tab's example arm also now shows the bench-calibrated defaults and calibrate\_arm.py's 2.5~s position steps.
    \item \textbf{Bug 17:} Fixed the 3D view drawing a run's arm links after the Spacecraft Designer's Chaser had swapped the arm for a docking cone. A run's arm now shows only while the Chaser carries it; otherwise the page asks whether that was intended and whether to show it anyway (with its mount where the run had it), not now, or never again.
    \item \textbf{Bug 18:} Fixed the Chaser's ZED2 model being added again on every change to the design, which left a stray copy in the 3D view once Use CAD was switched off.
    \item \textbf{Bug 19:} Fixed the Vision tab showing a camera's old settings after a change, for the DVXplorer and USB cameras (and the ZED's saved settings as it opened): SPOT Vision listed the settings before applying the change.
\end{itemize}

\section*{\textcolor{primaryColor}{Ongoing Issues}}
\begin{itemize}[label={\textcolor{primaryColor}{\textbullet}}]
    \item \textbf{Issue 1:} The arm calibration has now run on the real arm, and found that the wrist servo (ID 3) measures its current wrongly, and differently in each direction. Under torque control it would push harder one way than the other, so it needs replacing before the arm is calibrated or flown. Until then the friction values remain placeholders, and the controller gains may need retuning afterwards.
    \item \textbf{Issue 2:} Nothing has been tested on Windows yet, including starting the console and installing Python from SPOT.
    \item \textbf{Issue 3:} The Spacecraft Designer is a first version: the attachments don't affect the mass properties or the 2D figures yet, and the Experiment tab draws the project's design rather than the lab console's.
    \item \textbf{Issue 4:} Interactive Mode flies each platform straight at its goal, without any knowledge of the others, so their paths can still cross. Keep a hand near Esc.
    \item \textbf{Issue 5:} Wi-Fi power saving is still enabled on the platforms, which can delay their messages to the lab console.
    \item \textbf{Issue 6:} The table slope, and the Obstacle's thrusters and thrust decay, were calibrated in the lab on October 14th. The Chaser's and Target's thrusters have not been measured yet, so they use the nominal thrust and the Obstacle's decay.
    \item \textbf{Issue 7:} A few features of the old GUI have not been brought across yet: emptying the platforms' MATLAB\_ws and ExperimentSoftware folders, deploying to one platform at a time, video export, and custom drawings.
\end{itemize}

\section*{\textcolor{primaryColor}{Acknowledgments}}
Special thanks to the following people who assisted with the above software changes:
% \begin{itemize}
%     \item \textbf{Name}: Thank you for ...
% \end{itemize}

\end{document}
